\documentclass[11pt]{article}

\usepackage{amsmath,amssymb,amsthm}
\usepackage{geometry}
\usepackage[hidelinks]{hyperref}
\newtheorem{theorem}{Theorem}[section]
\newtheorem{proposition}[theorem]{Proposition}
\newtheorem{lemma}[theorem]{Lemma}
\newtheorem{corollary}[theorem]{Corollary}

\theoremstyle{definition}
\newtheorem{definition}[theorem]{Definition}

\theoremstyle{remark}
\newtheorem*{remark}{Remark}
\geometry{margin=1in}

\title{Benford-Weighted Digit Sums and Second-Order Asymptotics for Exponential Sequences}
\author{Radhe Dhoundiyal \\ \small{Independent Researcher}}
\date{September 5, 2026}

\begin{document}

\maketitle

\begin{abstract}
Let \(b\ge2\) be an integer base, and let \(M_L\) denote the
\(L\)-digit integer significand induced by the logarithmic, or Benford,
mantissa measure. Starting from the classical prefix law
\[
\mathbb P(M_L=m)=\log_b\!\left(\frac{m+1}{m}\right),
\]
we derive an exact finite evaluation for the Benford-weighted
base-\(b\) digit sum \(\mathbb E[s_b(M_L)]\). This yields an exact
finite formula for the cumulative signed deviation of significant-digit
expectations from the uniform digit mean and, in the limit,
\[
c_b=
\frac b2-\frac{b-1}{2}\log_b(2\pi).
\]
We also obtain a two-term finite-length expansion with an explicit
Stirling-remainder bound.

We then transfer this constant to deterministic exponential sequences.
For multiplicatively independent integers \(a,b\ge2\), a finite
discrepancy estimate shows that the mean sum of the first \(J\)
significant digits of \(a^n\) equals
\[
J\frac{b-1}{2}-c_b
\]
up to explicit errors. Combining this with classical quantitative
discrepancy bounds for \(\{n\log_b a\}\) gives, for a logarithmically
growing leading window \(J_N\),
\[
\mathcal S_{N,J_N}(a;b)
=
J_N\frac{b-1}{2}-c_b+o(1).
\]
Thus the same cumulative Benford bias appears as a bounded
second-order correction in a deterministic exponential sequence.
\end{abstract}

% ============================================================
% Section 1
% ============================================================

\section{Introduction and main results}
\label{sec:introduction}

The significant digits of exponential sequences provide a classical
meeting point of uniform distribution theory, Benford's law, and digital
number theory.  If \(b\ge 2\) is an integer base and
\(\Theta\) is uniformly distributed on \([0,1)\), then the random
mantissa
\[
    X=b^\Theta\in[1,b)
\]
induces the logarithmic, or Benford, distribution on significant
digits.  Equivalently, questions about the significant digits of
deterministic powers \(a^n\) are naturally related to the rotation
\[
    \{n\log_b a\},
\]
whenever \(a\) and \(b\) are multiplicatively independent.

The distribution of a fixed significant prefix in the logarithmic
mantissa model is classical.  In particular, if \(M_L\) denotes the
integer formed by the first \(L\) significant base-\(b\) digits of
\(b^\Theta\), then
\begin{equation}
    \mathbb P(M_L=m)
    =
    \log_b\!\left(\frac{m+1}{m}\right),
    \qquad
    b^{L-1}\le m\le b^L-1.
    \label{eq:classical-prefix-law-intro}
\end{equation}
The corresponding marginal laws of individual significant digits,
their relation to uniform distribution modulo one, and the convergence
of deeper significant digits toward the uniform digit distribution are
also well established; see, among others,
\cite{Newcomb1881,Benford1938,Diaconis1977,Hill1995Phenomenon,
Hill1995Statistical,BergerHill2011,BergerHill2015}.
Closed Gamma-function expressions for the individual higher-digit
marginals were obtained by Balado and Silvestre
\cite{BaladoSilvestre2021,BaladoSilvestre2024}.  We do not claim
novelty for these fixed-position distribution laws.

The purpose of this paper is instead to study an aggregate quantity that
retains the systematic bias of the early significant positions.

Let \(D_j^{(b)}\) denote the \(j\)-th significant base-\(b\) digit
under the logarithmic mantissa model, and set
\[
    \mu_b=\frac{b-1}{2},
\]
the mean of a uniformly distributed digit in
\(\{0,\ldots,b-1\}\).
For \(L\ge1\), define the finite cumulative signed significant-digit
bias by
\begin{equation}
    c_{b,L}
    =
    \sum_{j=1}^{L}
    \left(
        \mu_b-\mathbb E[D_j^{(b)}]
    \right).
    \label{eq:cbl-definition-intro}
\end{equation}
The adjective ``signed'' is important: the quantity need not be
positive for every base.

Our first main result gives an exact finite evaluation of the
corresponding Benford-weighted digit sum.  Let \(s_b(m)\) denote the
ordinary sum of the base-\(b\) digits of the integer \(m\).

\begin{theorem}[Exact Benford-weighted digit sum]
\label{thm:intro-finite-digit-sum}
For every integer \(b\ge2\) and every \(L\ge1\),
\begin{equation}
\boxed{
    \mathbb E[s_b(M_L)]
    =
    b^L-b+L(b-1)
    -
    \log_b
    \left(
        \frac{(b^L-1)!}
        {\bigl((b^{L-1}-1)!\bigr)^b}
    \right).
}
\label{eq:intro-finite-digit-sum}
\end{equation}
Equivalently,
\begin{equation}
\boxed{
    c_{b,L}
    =
    \frac{
        \ln\Gamma(b^L)
        -
        b\ln\Gamma(b^{L-1})
    }{\ln b}
    -
    (b^L-b)
    -
    \frac{L(b-1)}{2}.
}
\label{eq:intro-cbl-gamma}
\end{equation}
\end{theorem}

The proof uses the classical prefix law
\eqref{eq:classical-prefix-law-intro} together with the elementary
base-\(b\) recursion
\[
    s_b(bq+r)=s_b(q)+r,
    \qquad 0\le r\le b-1.
\]
Although the probability law in
\eqref{eq:classical-prefix-law-intro} is classical, we have not found
the weighted evaluation
\eqref{eq:intro-finite-digit-sum}, or the equivalent finite cumulative
identity \eqref{eq:intro-cbl-gamma}, in the literature surveyed.

Passing to the limit produces a simple closed constant.

\begin{theorem}[Cumulative significant-digit bias]
\label{thm:intro-cb}
For every integer base \(b\ge2\), the series
\[
    c_b
    =
    \sum_{j=1}^{\infty}
    \left(
        \frac{b-1}{2}
        -
        \mathbb E[D_j^{(b)}]
    \right)
\]
converges absolutely and satisfies
\begin{equation}
\boxed{
    c_b
    =
    \frac{b}{2}
    -
    \frac{b-1}{2}\log_b(2\pi).
}
\label{eq:intro-cb}
\end{equation}
\end{theorem}

For example, in base \(10\),
\begin{equation}
    c_{10}
    =
    5-\frac{9}{2}\log_{10}(2\pi)
    =
    1.408190592388482\ldots.
    \label{eq:intro-c10}
\end{equation}
Thus the cumulative effect of the biased early significant positions
is governed by the same \(2\pi\) constant that appears in Stirling's
formula.

The finite identity also gives precise control of the approach to the
limit.

\begin{theorem}[Finite-length expansion]
\label{thm:intro-tail}
For every \(b\ge2\) and \(L\ge1\),
\begin{equation}
\begin{aligned}
    c_b-c_{b,L}
    &=
    \frac{b^2-1}{12\ln b}\,b^{-L}
    -
    \frac{b^4-1}{360\ln b}\,b^{-3L}
    +
    \varepsilon_{b,L},
\end{aligned}
\label{eq:intro-tail-expansion}
\end{equation}
where
\begin{equation}
\boxed{
    |\varepsilon_{b,L}|
    \le
    \frac{b^6+1}{1260\ln b}\,b^{-5L}.
}
\label{eq:intro-tail-remainder}
\end{equation}
\end{theorem}

The individual significant-digit distributions admit corresponding
deep-position expansions.  These are useful below, but are not the
principal novelty of the paper.  Expanding the known exact
higher-significant-digit law yields, for \(j\ge2\) and
\(r\in\{0,\ldots,b-1\}\),
\begin{equation}
    \mathbb P(D_j^{(b)}=r)
    =
    \frac1b
    +
    \frac{(b-1)(b-1-2r)}
         {2b\ln b}\,b^{-j}
    +
    O_b(b^{-2j}),
    \label{eq:intro-digit-expansion}
\end{equation}
and consequently
\begin{equation}
    \mathbb E[D_j^{(b)}]
    =
    \frac{b-1}{2}
    -
    \frac{(b-1)^2(b+1)}
         {12\ln b}\,b^{-j}
    +
    O_b(b^{-2j}).
    \label{eq:intro-mean-digit-expansion}
\end{equation}
These formulas should be viewed as quantitative consequences of the
known higher-digit distribution rather than as independent distribution
theorems.

\subsection{From the logarithmic model to exponential sequences}
\label{subsec:intro-deterministic}

The second aim of the paper is to determine whether the constant
\(c_b\), obtained in the logarithmic mantissa model, persists in an
actual deterministic exponential sequence.

Let \(a,b\ge2\) be multiplicatively independent integers and write
\[
    \alpha=\log_b a.
\]
Then \(\alpha\) is irrational.  For \(n\ge1\), let
\[
    L_n=\lfloor n\alpha\rfloor+1
\]
be the number of base-\(b\) digits of \(a^n\), and let
\(D_j^{(b)}(a^n)\) denote its \(j\)-th significant digit.

For \(N,J\ge1\), define the leading-window digit-sum average
\begin{equation}
    \mathcal S_{N,J}(a;b)
    =
    \frac1N
    \sum_{\substack{1\le n\le N\\L_n\ge J}}
    \sum_{j=1}^{J}D_j^{(b)}(a^n).
    \label{eq:intro-SNJ}
\end{equation}
The normalization is by \(N\); powers with fewer than \(J\) digits
contribute nothing to the sum.

Let \(D_N(\alpha)\) denote the ordinary interval discrepancy of
\[
    \{\alpha\},\{2\alpha\},\ldots,\{N\alpha\}.
\]
That is,
\[
    D_N(\alpha)
    =
    \sup_{0\le u<v\le1}
    \left|
        \frac1N
        \#\{1\le n\le N:\{n\alpha\}\in[u,v)\}
        -(v-u)
    \right|.
\]

The digit cylinders associated with the first \(J\) significant
positions have exponentially increasing combinatorial complexity, but
ordinary discrepancy still gives an explicit finite transfer estimate.

\begin{theorem}[Finite deterministic transfer]
\label{thm:intro-deterministic-finite}
Let \(a,b\ge2\) be multiplicatively independent.  Let
\[
    e_J(N)
    =
    \#\{1\le n\le N:L_n<J\}.
\]
Then
\begin{equation}
\boxed{
\begin{aligned}
    \left|
        \mathcal S_{N,J}(a;b)
        -
        \left(
            J\frac{b-1}{2}-c_{b,J}
        \right)
    \right|
    \le{}&
    \frac{b-1}{2}\,b^J D_N(\alpha)
    \\
    &+
    \frac{(b-1)J\,e_J(N)}{N}.
\end{aligned}
}
\label{eq:intro-finite-transfer}
\end{equation}
\end{theorem}

Since
\[
    e_J(N)=O_{a,b}(J),
\]
Theorem~\ref{thm:intro-deterministic-finite} and
Theorem~\ref{thm:intro-tail} imply the following second-order
asymptotic.

\begin{theorem}[Second-order leading-window asymptotic]
\label{thm:intro-second-order}
Let \(J=J(N)\) satisfy
\[
    J(N)\longrightarrow\infty,
    \qquad
    b^{J(N)}D_N(\log_b a)\longrightarrow0,
    \qquad
    \frac{J(N)^2}{N}\longrightarrow0.
\]
Then
\begin{equation}
\boxed{
    \mathcal S_{N,J(N)}(a;b)
    =
    J(N)\frac{b-1}{2}
    -
    c_b
    +
    o(1).
}
\label{eq:intro-second-order}
\end{equation}
More quantitatively,
\begin{equation}
\begin{aligned}
    \mathcal S_{N,J}(a;b)
    =
    J\frac{b-1}{2}
    -
    c_b
    &+
    O_b(b^{-J})
    +
    O_b\!\left(b^J D_N(\log_b a)\right)
    \\
    &+
    O_{a,b}\!\left(\frac{J^2}{N}\right).
\end{aligned}
\label{eq:intro-second-order-quantitative}
\end{equation}
\end{theorem}

The point of
\eqref{eq:intro-second-order} is not the first-order term
\(J(b-1)/2\): first-order digit and block statistics in growing regions
of exponential sequences have been studied previously by discrepancy
and exponential-sum methods, notably by Fuchs
\cite{Fuchs2002} and the literature cited there.
The distinction here is that the accumulated nonuniformity of the
significant positions is retained on the absolute \(O(1)\) scale,
and its limiting contribution is exactly the same constant \(c_b\)
that arises in the logarithmic mantissa model.

Classical quantitative discrepancy estimates for
\(\alpha=\log_b a\), based ultimately on Diophantine estimates for
linear forms in logarithms, yield a power saving for every fixed
multiplicatively independent pair \((a,b)\); see
\cite{Fuchs2002,KissTichy1989I,KissTichy1989II,BakerWustholz,DrmotaTichy}.
Consequently, one obtains logarithmically growing windows as a direct
corollary.

\begin{corollary}[Logarithmically growing deterministic window]
\label{cor:intro-log-window}
Fix multiplicatively independent integers \(a,b\ge2\).
Suppose
\[
    D_N(\log_b a)\ll_{a,b}N^{-\delta_{a,b}}
\]
for some \(\delta_{a,b}>0\), as supplied by the classical quantitative
discrepancy theory cited above.  For any
\[
    0<\kappa<\delta_{a,b},
    \qquad
    J_N=\left\lfloor\kappa\log_b N\right\rfloor,
\]
one has
\begin{equation}
\boxed{
    \mathcal S_{N,J_N}(a;b)
    =
    J_N\frac{b-1}{2}
    -
    c_b
    +
    o(1).
}
\label{eq:intro-log-window}
\end{equation}
\end{corollary}

Thus the constant \(c_b\) has two interpretations.  It is first the
total signed positional bias of the significant digits under the
logarithmic mantissa measure, and it then reappears as a deterministic
boundary correction for a growing leading block of the powers \(a^n\).

\subsection{Relation to prior work}
\label{subsec:intro-prior-work}

The logarithmic law for leading significant digits goes back to
Newcomb and Benford
\cite{Newcomb1881,Benford1938}.  The relation between leading-digit
phenomena and uniform distribution modulo one was developed
systematically by Diaconis \cite{Diaconis1977}.  Hill's work
\cite{Hill1995Phenomenon,Hill1995Statistical} and the later treatment
of Berger and Hill \cite{BergerHill2011,BergerHill2015} provide a
modern framework for significant-digit distributions and Benford
measures.  Regularity of deeper significant digits has also been
studied in considerably greater generality
\cite{HillSchurger2005}.

Particularly relevant here is the work of Balado and Silvestre
\cite{BaladoSilvestre2021,BaladoSilvestre2024}, who give an exact
Gamma-function representation for the marginal distribution of the
\(j\)-th significant digit of a Benford variable.  That formula is
prior work and is not claimed here.  The position-wise expansion
\eqref{eq:intro-digit-expansion} may be obtained by expanding this
known exact distribution.

The ordinary sum-of-digits function has an extensive literature,
including the classical work of Trollope and Delange and later
Mellin-transform approaches
\cite{Trollope1968,Delange1975,FlajoletEtAl1994}.
Those results primarily concern counting or summation over integers
with uniform arithmetic weights.  The finite expectation considered
in Theorem~\ref{thm:intro-finite-digit-sum} instead weights the
\(L\)-digit integer significands according to the logarithmic masses
\[
    \log_b\frac{m+1}{m}.
\]
To the best of our knowledge, the resulting exact factorial identity
and the closed evaluation \eqref{eq:intro-cb} have not previously been
recorded.

For deterministic exponential sequences, Fuchs
\cite{Fuchs2002} established quantitative first-order results for
digital blocks in growing regions using discrepancy and exponential
sums, building on earlier work in discrepancy and Diophantine
approximation.  We therefore do not claim novelty for the use of
discrepancy, for first-order uniform digit frequencies, or for the
existence of logarithmically growing controlled windows.  Our
second-order statement retains the bounded contribution left after
subtracting the uniform first-order digit mean.  We have not found the
explicit constant \(-c_b\), or an equivalent evaluation, in the
earlier exponential-sequence literature surveyed.

Accordingly, the principal contribution of the present paper is the
chain
\[
\boxed{
\begin{gathered}
    \text{classical Benford prefix law}
    \\
    \Downarrow
    \\
    \text{exact Benford-weighted finite digit-sum identity}
    \\
    \Downarrow
    \\
    c_b=
    \frac b2-\frac{b-1}{2}\log_b(2\pi)
    \\
    \Downarrow
    \\
    \text{the same }c_b
    \text{ as a deterministic second-order correction.}
\end{gathered}
}
\]

\subsection{Scope and organization}
\label{subsec:intro-organization}

The deterministic results in this paper concern a growing block of
\emph{leading significant digits}.  They do not establish an
asymptotic law for the complete base-\(b\) digit sum of \(a^n\), whose
number of digits is of order \(n\).  In particular, no independence,
mixing, central limit theorem, or variance law for the complete digit
sequence is assumed or proved here.

Section~\ref{sec:benford-preliminaries} recalls the logarithmic
significand model and the classical significant-digit laws needed
below.  Section~\ref{sec:finite-bias} proves the exact finite
Benford-weighted digit-sum identity and evaluates the cumulative
constant \(c_b\).  Section~\ref{sec:quantitative-asymptotics} develops
the finite-position and finite-length asymptotics.
Section~\ref{sec:deterministic-transfer} transfers the resulting
constant to deterministic exponential sequences by discrepancy
estimates. Section~\ref{sec:verification} records the independent computational
checks used to audit the formulas,
Section~\ref{sec:discussion} discusses the limitations of the method
and several open problems, and
Section~\ref{sec:conclusion} concludes the paper.

% ============================================================
% Section 2
% ============================================================

\section{Benford prefixes and significant digits}
\label{sec:benford-preliminaries}

This section fixes notation and records the classical significant-digit
facts used later.  None of the distribution laws in this section are
claimed as new.  The finite prefix law is standard Benford theory, and
the Gamma-function form of the higher significant-digit marginals is
due to Balado and Silvestre
\cite{BaladoSilvestre2021,BaladoSilvestre2024}.

Throughout, \(b\ge2\) is an integer base and
\[
    \Theta\sim\operatorname{Unif}[0,1).
\]
We write
\[
    X=b^\Theta\in[1,b).
\]
Equivalently, \(X\) has density
\[
    \frac{1}{x\ln b},
    \qquad 1\le x<b,
\]
with respect to Lebesgue measure.

For \(j\ge1\), let \(D_j^{(b)}\) denote the \(j\)-th significant
base-\(b\) digit of \(X\).  Thus
\[
    D_1^{(b)}\in\{1,\ldots,b-1\},
\]
whereas
\[
    D_j^{(b)}\in\{0,\ldots,b-1\},
    \qquad j\ge2.
\]

It will also be convenient to use the ordinary base-\(b\) digit-sum
function
\[
    s_b(m),
\]
defined for nonnegative integers \(m\).

\subsection{Finite integer significands}
\label{subsec:finite-significands}

For \(L\ge1\), define the \(L\)-digit integer significand
\begin{equation}
    M_L
    =
    \left\lfloor b^{L-1+\Theta}\right\rfloor.
    \label{eq:ML-definition}
\end{equation}
Then
\[
    b^{L-1}\le M_L\le b^L-1,
\]
and the base-\(b\) digits of \(M_L\) are precisely
\[
    D_1^{(b)},D_2^{(b)},\ldots,D_L^{(b)}.
\]
Consequently,
\begin{equation}
    s_b(M_L)
    =
    \sum_{j=1}^{L}D_j^{(b)}.
    \label{eq:digit-sum-prefix}
\end{equation}

The distribution of \(M_L\) is the standard Benford distribution on
\(L\)-digit integer significands.

\begin{proposition}[Classical Benford prefix law]
\label{prop:benford-prefix-law}
For
\[
    b^{L-1}\le m\le b^L-1,
\]
one has
\begin{equation}
    \mathbb P(M_L=m)
    =
    \log_b\!\left(\frac{m+1}{m}\right).
    \label{eq:benford-prefix-law}
\end{equation}
\end{proposition}

\begin{proof}
By definition,
\[
    M_L=m
\]
if and only if
\[
    m
    \le
    b^{L-1+\Theta}
    <
    m+1.
\]
Taking logarithms gives
\[
    \log_b\!\left(\frac{m}{b^{L-1}}\right)
    \le
    \Theta
    <
    \log_b\!\left(\frac{m+1}{b^{L-1}}\right).
\]
Since \(\Theta\) is uniform on \([0,1)\), the probability is the
length of this interval:
\[
    \log_b(m+1)-\log_b m
    =
    \log_b\!\left(\frac{m+1}{m}\right).
\]
\end{proof}

The proof is included only to establish notation and normalization.
Equation~\eqref{eq:benford-prefix-law} is classical; see, for example,
\cite{Diaconis1977,Hill1995Phenomenon,BergerHill2011,BergerHill2015}.

For \(L=1\), Proposition~\ref{prop:benford-prefix-law} yields the
usual first-digit law
\begin{equation}
    \mathbb P(D_1^{(b)}=d)
    =
    \log_b\!\left(1+\frac1d\right),
    \qquad
    d=1,\ldots,b-1.
    \label{eq:first-digit-law}
\end{equation}
A direct finite summation also gives
\begin{equation}
    \mathbb E[D_1^{(b)}]
    =
    (b-1)-\log_b((b-1)!).
    \label{eq:first-digit-expectation}
\end{equation}
We record \eqref{eq:first-digit-expectation} for later comparison;
it is an elementary consequence of the classical leading-digit law.

\subsection{Individual higher significant digits}
\label{subsec:higher-digit-laws}

Let \(j\ge2\) and
\[
    r\in\{0,\ldots,b-1\}.
\]
Put
\begin{equation}
    M=b^{j-2}.
    \label{eq:M-j-definition}
\end{equation}

A \(j\)-digit prefix whose final digit is \(r\) has the form
\[
    bq+r,
\]
where
\[
    M\le q\le bM-1.
\]
Summing the classical prefix probabilities therefore gives the
\(j\)-th digit marginal.

\begin{proposition}[Classical higher-digit marginal]
\label{prop:higher-digit-marginal}
For \(j\ge2\),
\begin{equation}
\boxed{
    \mathbb P(D_j^{(b)}=r)
    =
    \sum_{q=M}^{bM-1}
    \log_b
    \left(
        \frac{bq+r+1}{bq+r}
    \right).
}
\label{eq:higher-digit-log-sum}
\end{equation}
\end{proposition}

\begin{proof}
The events
\[
    \{M_j=bq+r\},
    \qquad M\le q\le bM-1,
\]
are disjoint and exhaust the possible \(j\)-digit prefixes having
final digit \(r\).  Applying
Proposition~\ref{prop:benford-prefix-law} with \(L=j\) and summing the
corresponding probabilities gives
\eqref{eq:higher-digit-log-sum}.
\end{proof}

An equivalent closed form is obtained by converting the finite product
implicit in \eqref{eq:higher-digit-log-sum} to Gamma functions.
This representation is prior work of Balado and Silvestre
\cite{BaladoSilvestre2021,BaladoSilvestre2024}.

\begin{proposition}[Balado--Silvestre Gamma representation]
\label{prop:balado-silvestre-gamma}
For \(j\ge2\), \(M=b^{j-2}\), and
\(r\in\{0,\ldots,b-1\}\),
\begin{equation}
\boxed{
\begin{aligned}
    \mathbb P(D_j^{(b)}=r)
    &=
    \frac{1}{\ln b}
    \ln
    \left[
    \frac{
        \Gamma\!\left(bM+\frac{r+1}{b}\right)
        \Gamma\!\left(M+\frac{r}{b}\right)
    }{
        \Gamma\!\left(bM+\frac{r}{b}\right)
        \Gamma\!\left(M+\frac{r+1}{b}\right)
    }
    \right].
\end{aligned}
}
\label{eq:higher-digit-gamma}
\end{equation}
\end{proposition}

\begin{proof}
From \eqref{eq:higher-digit-log-sum},
\[
\begin{aligned}
    \mathbb P(D_j^{(b)}=r)
    &=
    \frac{1}{\ln b}
    \ln
    \prod_{q=M}^{bM-1}
    \frac{q+(r+1)/b}{q+r/b}.
\end{aligned}
\]
Using
\[
    \prod_{q=M}^{bM-1}(q+\lambda)
    =
    \frac{\Gamma(bM+\lambda)}
         {\Gamma(M+\lambda)}
\]
gives \eqref{eq:higher-digit-gamma}.
\end{proof}

\begin{remark}[Attribution]
\label{rem:gamma-attribution}
The derivation above is included for completeness and to fix the
present notation.  The Gamma representation
\eqref{eq:higher-digit-gamma} is not a contribution of this paper.
Balado and Silvestre gave this form in their 2021 conference paper and
in their later general treatment of number-representation elements
\cite{BaladoSilvestre2021,BaladoSilvestre2024}.
\end{remark}

The probabilities in
\eqref{eq:higher-digit-log-sum} satisfy
\[
    \sum_{r=0}^{b-1}
    \mathbb P(D_j^{(b)}=r)
    =
    1.
\]
Moreover,
\[
    \mathbb P(D_j^{(b)}=r)
    \longrightarrow
    \frac1b
    \qquad
    (j\to\infty),
\]
a classical manifestation of the asymptotic regularity of deeper
significant digits
\cite{Hill1995Phenomenon,HillSchurger2005,
BaladoSilvestre2024}.
A quantitative first correction to this limit will be recorded in
Section~\ref{sec:quantitative-asymptotics}, where it is used as
supporting information rather than as a principal novelty claim.

\subsection{From position-wise laws to a finite digit sum}
\label{subsec:position-to-sum}

By \eqref{eq:digit-sum-prefix} and linearity of expectation,
\begin{equation}
    \mathbb E[s_b(M_L)]
    =
    \sum_{j=1}^{L}\mathbb E[D_j^{(b)}].
    \label{eq:expectation-sum-positions}
\end{equation}
Thus one could, in principle, study the cumulative effect of the
nonuniform significant positions by summing the individual marginal
expectations.

For the main results, however, it is more effective to work directly
with the finite random integer \(M_L\).  Combining
Proposition~\ref{prop:benford-prefix-law} with
\eqref{eq:digit-sum-prefix} gives
\begin{equation}
    \mathbb E[s_b(M_L)]
    =
    \sum_{m=b^{L-1}}^{b^L-1}
    s_b(m)
    \log_b\!\left(\frac{m+1}{m}\right).
    \label{eq:weighted-digit-sum-start}
\end{equation}

Equation~\eqref{eq:weighted-digit-sum-start} is the starting point of
the exact finite evaluation in the next section.  The probability
weights are classical; the issue is whether the logarithmically
weighted digital sum itself admits a useful closed form.
Section~\ref{sec:finite-bias} shows that it does.

% ============================================================
% Section 3
% ============================================================

\section{Exact Benford-weighted digit sums and cumulative bias}
\label{sec:finite-bias}

We now evaluate the logarithmically weighted digit sum introduced in
\eqref{eq:weighted-digit-sum-start}.  The input distribution is the
classical Benford prefix law; the new ingredient is that the weighted
sum collapses to a finite factorial expression.

For convenience, write
\begin{equation}
    E_{b,L}
    =
    \mathbb E[s_b(M_L)].
    \label{eq:E-bL-definition}
\end{equation}
By Proposition~\ref{prop:benford-prefix-law},
\begin{equation}
    E_{b,L}
    =
    \sum_{m=b^{L-1}}^{b^L-1}
    s_b(m)
    \log_b\!\left(\frac{m+1}{m}\right).
    \label{eq:E-bL-raw}
\end{equation}

The elementary identity
\begin{equation}
    s_b(bq+r)=s_b(q)+r,
    \qquad
    0\le r\le b-1,
    \label{eq:digit-sum-recursion}
\end{equation}
allows the sum to be reduced recursively.

\subsection{A finite block identity}
\label{subsec:finite-block-identity}

We first isolate the contribution of the final base-\(b\) digit.

\begin{lemma}
\label{lem:last-digit-block}
Let \(L\ge2\), and put
\[
    A=b^{L-2},
    \qquad
    B=b^{L-1}=bA.
\]
Then
\begin{equation}
    E_{b,L}
    =
    E_{b,L-1}
    +
    R_{b,L},
    \label{eq:E-recursion}
\end{equation}
where
\begin{equation}
\boxed{
    R_{b,L}
    =
    (b-1)(B+1)
    -
    \log_b
    \left(
        \frac{
            (bB-1)!\,(A-1)!^b
        }{
            (B-1)!^{\,b+1}
        }
    \right).
}
\label{eq:R-bL-closed}
\end{equation}
\end{lemma}

\begin{proof}
Every integer
\[
    m\in[B,bB-1]
\]
can be written uniquely as
\[
    m=bq+r,
\]
with
\[
    A\le q\le B-1,
    \qquad
    0\le r\le b-1.
\]
Using \eqref{eq:digit-sum-recursion} in
\eqref{eq:E-bL-raw} gives
\begin{align}
    E_{b,L}
    &=
    \sum_{q=A}^{B-1}
    \sum_{r=0}^{b-1}
    \bigl(s_b(q)+r\bigr)
    \log_b
    \left(
        \frac{bq+r+1}{bq+r}
    \right)
    \notag\\
    &=
    \sum_{q=A}^{B-1}
    s_b(q)
    \sum_{r=0}^{b-1}
    \log_b
    \left(
        \frac{bq+r+1}{bq+r}
    \right)
    +
    R_{b,L},
    \label{eq:block-split}
\end{align}
where
\begin{equation}
    R_{b,L}
    =
    \sum_{q=A}^{B-1}
    \sum_{r=0}^{b-1}
    r
    \log_b
    \left(
        \frac{bq+r+1}{bq+r}
    \right).
    \label{eq:R-bL-definition}
\end{equation}

The inner logarithmic sum in the first term telescopes:
\[
    \sum_{r=0}^{b-1}
    \log_b
    \left(
        \frac{bq+r+1}{bq+r}
    \right)
    =
    \log_b
    \left(
        \frac{q+1}{q}
    \right).
\]
Hence the first term in \eqref{eq:block-split} is precisely
\[
    E_{b,L-1}.
\]

It remains to evaluate \(R_{b,L}\).
For fixed \(q\), discrete summation by parts gives
\begin{align}
    &\sum_{r=0}^{b-1}
    r
    \log_b
    \left(
        \frac{bq+r+1}{bq+r}
    \right)
    \notag\\
    &\qquad=
    \log_b
    \left(
        \frac{
            (bq+b)^{b-1}
        }{
            \prod_{r=1}^{b-1}(bq+r)
        }
    \right).
    \label{eq:fixed-q-last-digit}
\end{align}

Multiplying \eqref{eq:fixed-q-last-digit} over
\(A\le q\le B-1\) yields
\begin{equation}
    R_{b,L}
    =
    \log_b
    \left[
        \frac{
            \displaystyle
            \prod_{q=A}^{B-1}
            (b(q+1))^{b-1}
        }{
            \displaystyle
            \prod_{q=A}^{B-1}
            \prod_{r=1}^{b-1}(bq+r)
        }
    \right].
    \label{eq:R-product}
\end{equation}

The numerator is
\begin{equation}
    b^{(b-1)(B-A)}
    \left(
        \frac{B!}{A!}
    \right)^{b-1}.
    \label{eq:R-numerator}
\end{equation}
The denominator consists of all integers from \(B+1\) to \(bB-1\)
that are not divisible by \(b\).  Therefore
\begin{align}
    \prod_{q=A}^{B-1}
    \prod_{r=1}^{b-1}(bq+r)
    &=
    \frac{
        (bB-1)!/B!
    }{
        \displaystyle
        \prod_{k=A+1}^{B-1}bk
    }
    \notag\\
    &=
    \frac{
        (bB-1)!\,A!
    }{
        B!\,
        b^{B-A-1}
        (B-1)!
    }.
    \label{eq:R-denominator}
\end{align}

Substituting
\eqref{eq:R-numerator} and
\eqref{eq:R-denominator} into
\eqref{eq:R-product}, and using
\[
    B=bA,
\]
one obtains
\begin{equation}
    \frac{
        \displaystyle
        \prod_{q=A}^{B-1}
        (b(q+1))^{b-1}
    }{
        \displaystyle
        \prod_{q=A}^{B-1}
        \prod_{r=1}^{b-1}(bq+r)
    }
    =
    b^{(b-1)(B+1)}
    \frac{
        (B-1)!^{\,b+1}
    }{
        (A-1)!^b\,(bB-1)!
    }.
    \label{eq:R-product-simplified}
\end{equation}
Taking \(\log_b\) gives
\eqref{eq:R-bL-closed}.
\end{proof}


\subsection{Exact finite Benford-weighted digit sum}
\label{subsec:exact-finite-digit-sum}

We can now prove the first main theorem stated in the introduction.

\begin{proof}[Proof of Theorem~\ref{thm:intro-finite-digit-sum}]
We argue by induction on \(L\).

For \(L=1\),
\[
    M_1=D_1^{(b)},
\]
and hence
\begin{align}
    E_{b,1}
    &=
    \sum_{d=1}^{b-1}
    d\log_b\!\left(\frac{d+1}{d}\right)
    \notag\\
    &=
    (b-1)-\log_b((b-1)!).
    \label{eq:E-base-case}
\end{align}
This agrees with
\eqref{eq:intro-finite-digit-sum}, since \(0!=1\).

Now suppose \(L\ge2\) and that
\begin{equation}
    E_{b,L-1}
    =
    b^{L-1}-b+(L-1)(b-1)
    -
    \log_b
    \left(
        \frac{
            (b^{L-1}-1)!
        }{
            ((b^{L-2}-1)!)^b
        }
    \right).
    \label{eq:E-induction-hypothesis}
\end{equation}
Again put
\[
    A=b^{L-2},
    \qquad
    B=b^{L-1}.
\]
By Lemma~\ref{lem:last-digit-block},
\begin{align}
    E_{b,L}
    &=
    E_{b,L-1}
    +
    (b-1)(B+1)
    \notag\\
    &\qquad
    -
    \log_b
    \left(
        \frac{
            (bB-1)!\,(A-1)!^b
        }{
            (B-1)!^{\,b+1}
        }
    \right).
    \label{eq:E-induction-step-start}
\end{align}

The elementary terms satisfy
\[
    B-b+(L-1)(b-1)+(b-1)(B+1)
    =
    bB-b+L(b-1).
\]
Since \(bB=b^L\), this is
\[
    b^L-b+L(b-1).
\]

The logarithmic terms combine as
\begin{align}
    &
    \frac{(B-1)!}{(A-1)!^b}
    \cdot
    \frac{(bB-1)!\,(A-1)!^b}
         {(B-1)!^{\,b+1}}
    \notag\\
    &\qquad=
    \frac{(bB-1)!}{(B-1)!^b}
    =
    \frac{
        (b^L-1)!
    }{
        ((b^{L-1}-1)!)^b
    }.
    \label{eq:E-log-combination}
\end{align}
Therefore
\[
\boxed{
    E_{b,L}
    =
    b^L-b+L(b-1)
    -
    \log_b
    \left(
        \frac{
            (b^L-1)!
        }{
            ((b^{L-1}-1)!)^b
        }
    \right),
}
\]
which completes the induction.
\end{proof}

\begin{remark}
\label{rem:finite-identity-novelty}
The probability measure used in
Theorem~\ref{thm:intro-finite-digit-sum} is not new:
\[
    \mathbb P(M_L=m)
    =
    \log_b\!\left(\frac{m+1}{m}\right)
\]
is the classical Benford law for the \(L\)-digit integer significand.
The content of the theorem is the exact evaluation obtained after
inserting the digital statistic \(s_b(m)\) into this logarithmically
weighted measure.
\end{remark}


\subsection{Finite cumulative significant-digit bias}
\label{subsec:finite-cumulative-bias}

Recall that
\[
    \mu_b=\frac{b-1}{2}
\]
and
\[
    c_{b,L}
    =
    \sum_{j=1}^{L}
    \left(
        \mu_b-\mathbb E[D_j^{(b)}]
    \right).
\]
Since
\[
    s_b(M_L)
    =
    \sum_{j=1}^{L}D_j^{(b)},
\]
linearity of expectation gives
\begin{equation}
    c_{b,L}
    =
    L\frac{b-1}{2}
    -
    E_{b,L}.
    \label{eq:cbl-from-E}
\end{equation}

Substituting
Theorem~\ref{thm:intro-finite-digit-sum} into
\eqref{eq:cbl-from-E} gives
\begin{align}
    c_{b,L}
    &=
    \log_b
    \left(
        \frac{
            (b^L-1)!
        }{
            ((b^{L-1}-1)!)^b
        }
    \right)
    -
    (b^L-b)
    -
    \frac{L(b-1)}{2}.
    \label{eq:cbl-factorial}
\end{align}
Using
\[
    n!=\Gamma(n+1)
\]
yields the equivalent Gamma representation
\begin{equation}
\boxed{
    c_{b,L}
    =
    \frac{
        \ln\Gamma(b^L)
        -
        b\ln\Gamma(b^{L-1})
    }{\ln b}
    -
    (b^L-b)
    -
    \frac{L(b-1)}{2}.
}
\label{eq:cbl-gamma-main}
\end{equation}

Thus the cumulative effect of the first \(L\) significant positions
has an exact finite form even though the individual digit
distributions are not uniform.


\subsection{The limiting bias constant}
\label{subsec:limiting-bias}

We next evaluate the limit of
\eqref{eq:cbl-gamma-main}.

\begin{proof}[Proof of Theorem~\ref{thm:intro-cb}]
Set
\[
    x=b^{L-1},
\]
so that
\[
    b^L=bx.
\]
The classical Stirling expansion gives
\begin{equation}
    \ln\Gamma(x)
    =
    \left(x-\frac12\right)\ln x
    -
    x
    +
    \frac12\ln(2\pi)
    +
    O(x^{-1})
    \qquad
    (x\to\infty).
    \label{eq:stirling-leading}
\end{equation}

Applying \eqref{eq:stirling-leading} at \(x\) and \(bx\) gives
\begin{align}
    \ln\Gamma(bx)-b\ln\Gamma(x)
    &=
    \left(bx-\frac12\right)\ln b
    +
    \frac{b-1}{2}\ln x
    \notag\\
    &\qquad
    -
    \frac{b-1}{2}\ln(2\pi)
    +
    O(x^{-1}).
    \label{eq:gamma-difference-stirling}
\end{align}
Since
\[
    \ln x=(L-1)\ln b,
\]
division by \(\ln b\) yields
\begin{align}
    \frac{
        \ln\Gamma(b^L)
        -
        b\ln\Gamma(b^{L-1})
    }{\ln b}
    &=
    b^L
    -
    \frac{b}{2}
    +
    \frac{L(b-1)}{2}
    \notag\\
    &\qquad
    -
    \frac{b-1}{2}\log_b(2\pi)
    +
    O_b(b^{-L}).
    \label{eq:gamma-difference-final}
\end{align}

Substituting this into
\eqref{eq:cbl-gamma-main}, all terms depending linearly on
\(b^L\) and \(L\) cancel.  We obtain
\begin{equation}
    c_{b,L}
    =
    \frac{b}{2}
    -
    \frac{b-1}{2}\log_b(2\pi)
    +
    O_b(b^{-L}).
    \label{eq:cbl-leading-limit}
\end{equation}
Therefore
\begin{equation}
\boxed{
    c_b
    =
    \lim_{L\to\infty}c_{b,L}
    =
    \frac{b}{2}
    -
    \frac{b-1}{2}\log_b(2\pi).
}
\label{eq:cb-main}
\end{equation}

The absolute convergence of
\[
    \sum_{j\ge1}
    \left(
        \frac{b-1}{2}
        -
        \mathbb E[D_j^{(b)}]
    \right)
\]
follows from the quantitative higher-position estimate proved in
Section~\ref{sec:quantitative-asymptotics}.
\end{proof}

\begin{remark}[Signed nature of the bias]
\label{rem:signed-bias}
The term ``bias'' in this paper is signed.
For example,
\[
    c_2<0,
    \qquad
    c_3<0,
\]
whereas
\[
    c_{10}
    =
    5-\frac92\log_{10}(2\pi)
    =
    1.408190592388482\ldots.
\]
Thus it would be misleading to call \(c_b\) a universal
``digit deficit'' across all bases.
\end{remark}


\subsection{Expected average of a finite significant prefix}
\label{subsec:model-average}

Define
\begin{equation}
    A_L^{(b)}
    =
    \frac1L
    \sum_{j=1}^{L}D_j^{(b)}
    =
    \frac{s_b(M_L)}{L}.
    \label{eq:A-L-base-b}
\end{equation}
Then the definition of \(c_{b,L}\) gives the exact identity
\begin{equation}
\boxed{
    \mathbb E[A_L^{(b)}]
    =
    \frac{b-1}{2}
    -
    \frac{c_{b,L}}{L}.
}
\label{eq:model-average-exact}
\end{equation}
Consequently,
\begin{equation}
    \mathbb E[A_L^{(b)}]
    =
    \frac{b-1}{2}
    -
    \frac{c_b}{L}
    +
    O_b\!\left(\frac{b^{-L}}{L}\right),
    \label{eq:model-average-leading}
\end{equation}
with the sharper expansion and explicit remainder given in
Section~\ref{sec:quantitative-asymptotics}.

Equation~\eqref{eq:model-average-leading} provides the first
interpretation of \(c_b\): it is the finite boundary correction left
by the nonuniform early significant positions after the bulk digit
mean has approached the uniform value \((b-1)/2\).
In Section~\ref{sec:deterministic-transfer}, we show that the same
constant survives when the logarithmic mantissa measure is replaced
by a quantitative orbit average along the deterministic sequence
\(a^n\).

% ============================================================
% Section 4
% ============================================================

\section{Quantitative finite-position and finite-length asymptotics}
\label{sec:quantitative-asymptotics}

The preceding section evaluated the cumulative significant-digit bias
exactly.  We now record two quantitative consequences.

First, the known exact distribution of an individual higher significant
digit admits a simple first-order expansion about the uniform digit law.
Second, the exact finite formula for \(c_{b,L}\) yields a two-term
asymptotic expansion with an explicit Stirling-remainder bound.

The first result is included primarily as quantitative background for
the cumulative theory.  It should not be interpreted as a new
significant-digit distribution theorem: the exact marginal law and its
Gamma representation are prior work, as discussed in
Section~\ref{subsec:higher-digit-laws}.

\subsection{Asymptotics of an individual significant digit}
\label{subsec:individual-digit-asymptotics}

Let \(j\ge2\), and write
\[
    M=b^{j-2}.
\]
For
\[
    r\in\{0,\ldots,b-1\},
\]
Proposition~\ref{prop:balado-silvestre-gamma} gives
\begin{equation}
    \mathbb P(D_j^{(b)}=r)
    =
    \frac{1}{\ln b}
    \left[
        \Delta(bM)-\Delta(M)
    \right],
    \label{eq:digit-pmf-delta}
\end{equation}
where
\[
    \Delta(z)
    =
    \ln\Gamma(z+u)-\ln\Gamma(z+v),
\]
with
\[
    u=\frac{r+1}{b},
    \qquad
    v=\frac{r}{b}.
\]

We use the standard fixed-parameter Gamma-ratio expansion
\cite[\S5.11(iii)]{DLMFStirling}
\begin{equation}
    \ln\Gamma(z+u)-\ln\Gamma(z+v)
    =
    (u-v)\ln z
    +
    \frac{(u-v)(u+v-1)}{2z}
    +
    O_{u,v}(z^{-2})
    \label{eq:loggamma-ratio-first-order}
\end{equation}
as \(z\to\infty\).
Since \(b\) is fixed and \(r\) ranges over only finitely many values,
the remainder may be taken uniformly in
\(r\in\{0,\ldots,b-1\}\), with a constant depending only on \(b\).

\begin{proposition}[Deep-position expansion]
\label{prop:deep-position-expansion}
For fixed \(b\ge2\), uniformly in
\(r\in\{0,\ldots,b-1\}\),
\begin{equation}
\boxed{
    \mathbb P(D_j^{(b)}=r)
    =
    \frac1b
    +
    \frac{(b-1)(b-1-2r)}
         {2b\ln b}\,b^{-j}
    +
    O_b(b^{-2j})
}
\label{eq:deep-position-expansion}
\end{equation}
as \(j\to\infty\).
\end{proposition}

\begin{proof}
From
\[
    u-v=\frac1b
\]
and
\[
    u+v-1
    =
    \frac{2r+1-b}{b},
\]
equation~\eqref{eq:loggamma-ratio-first-order} gives
\begin{equation}
    \Delta(z)
    =
    \frac1b\ln z
    +
    \frac{2r+1-b}{2b^2z}
    +
    O_b(z^{-2}).
    \label{eq:delta-expansion}
\end{equation}
Applying this at \(z=bM\) and \(z=M\), we obtain
\begin{align}
    \Delta(bM)-\Delta(M)
    &=
    \frac1b\ln b
    \notag\\
    &\quad+
    \frac{2r+1-b}{2b^2}
    \left(
        \frac1{bM}-\frac1M
    \right)
    +
    O_b(M^{-2}).
    \label{eq:delta-difference}
\end{align}
Since
\[
    \frac1{bM}-\frac1M
    =
    -\frac{b-1}{bM},
\]
the correction term becomes
\[
    \frac{(b-1)(b-1-2r)}
         {2b^3M}.
\]
Using
\[
    M=b^{j-2},
    \qquad
    \frac{1}{b^3M}
    =
    \frac{1}{b}\,b^{-j},
\]
and dividing by \(\ln b\) in
\eqref{eq:digit-pmf-delta} yields
\eqref{eq:deep-position-expansion}.
Finally,
\[
    M^{-2}=b^{-2j+4}=O_b(b^{-2j}),
\]
which gives the stated remainder.
\end{proof}

\begin{remark}[Relation to prior work]
\label{rem:deep-position-prior-work}
The convergence
\[
    \mathbb P(D_j^{(b)}=r)\longrightarrow\frac1b
\]
is classical, and the exact Gamma representation from which
Proposition~\ref{prop:deep-position-expansion} is derived is due to
Balado and Silvestre
\cite{BaladoSilvestre2021,BaladoSilvestre2024}.
Thus
\eqref{eq:deep-position-expansion} is best viewed as an explicit
first correction obtained by expanding a known exact formula.
\end{remark}


\subsection{Expectation of a deep significant digit}
\label{subsec:mean-digit-asymptotics}

Summing the preceding expansion against the digit value gives the
corresponding mean.

\begin{corollary}
\label{cor:mean-digit-expansion}
For fixed \(b\ge2\),
\begin{equation}
\boxed{
    \mathbb E[D_j^{(b)}]
    =
    \frac{b-1}{2}
    -
    \frac{(b-1)^2(b+1)}
         {12\ln b}\,b^{-j}
    +
    O_b(b^{-2j})
}
\label{eq:mean-digit-expansion}
\end{equation}
as \(j\to\infty\).
\end{corollary}

\begin{proof}
By definition,
\[
    \mathbb E[D_j^{(b)}]
    =
    \sum_{r=0}^{b-1}
    r\,\mathbb P(D_j^{(b)}=r).
\]
The uniform term contributes
\[
    \frac1b
    \sum_{r=0}^{b-1}r
    =
    \frac{b-1}{2}.
\]
For the first correction, use
\begin{align}
    \sum_{r=0}^{b-1}
    r(b-1-2r)
    &=
    (b-1)\sum_{r=0}^{b-1}r
    -
    2\sum_{r=0}^{b-1}r^2
    \notag\\
    &=
    -\frac{b(b-1)(b+1)}{6}.
    \label{eq:r-correction-sum}
\end{align}
Multiplying by the coefficient
\[
    \frac{b-1}{2b\ln b}\,b^{-j}
\]
from
\eqref{eq:deep-position-expansion}
gives
\[
    -
    \frac{(b-1)^2(b+1)}
         {12\ln b}\,b^{-j}.
\]
The finite sum over \(r\) preserves the
\(O_b(b^{-2j})\) remainder.
\end{proof}

In particular,
\begin{equation}
    \frac{b-1}{2}
    -
    \mathbb E[D_j^{(b)}]
    =
    O_b(b^{-j}),
    \label{eq:bias-geometric-decay}
\end{equation}
and therefore
\[
    \sum_{j=1}^{\infty}
    \left|
        \frac{b-1}{2}
        -
        \mathbb E[D_j^{(b)}]
    \right|
    <
    \infty.
\]
This completes the absolute-convergence assertion in
Theorem~\ref{thm:intro-cb}.


\subsection{A two-term expansion for the cumulative tail}
\label{subsec:cumulative-tail-expansion}

We now return to the finite cumulative bias.
The exact formula
\eqref{eq:cbl-gamma-main} permits substantially sharper information
than the estimate
\[
    c_{b,L}=c_b+O_b(b^{-L})
\]
used in Section~\ref{subsec:limiting-bias}.

For \(x>0\), write Stirling's expansion in the form
\begin{equation}
\begin{aligned}
    \ln\Gamma(x)
    &=
    \left(x-\frac12\right)\ln x
    -
    x
    +
    \frac12\ln(2\pi)
    \\
    &\quad+
    \frac{1}{12x}
    -
    \frac{1}{360x^3}
    +
    \rho(x).
\end{aligned}
\label{eq:stirling-two-term}
\end{equation}
For positive real \(x\), the standard first-neglected-term remainder
bound for Stirling's series gives
\begin{equation}
    0<\rho(x)
    \le
    \frac{1}{1260x^5}.
    \label{eq:stirling-rho-bound}
\end{equation}
We use the positive-real remainder theorem in
\cite[\S5.11(i)--(ii)]{DLMFStirling}.

Define
\begin{equation}
    R(x)
    =
    \frac{1}{12x}
    -
    \frac{1}{360x^3}
    +
    \rho(x).
    \label{eq:R-stirling}
\end{equation}

\begin{proof}[Proof of Theorem~\ref{thm:intro-tail}]
Put
\[
    x=b^{L-1}.
\]
Then
\[
    bx=b^L.
\]
Substituting the Stirling expansion
\eqref{eq:stirling-two-term}
into the exact formula
\eqref{eq:cbl-gamma-main}, and subtracting from the limiting value
\eqref{eq:cb-main}, the elementary terms cancel exactly and give
\begin{equation}
    c_b-c_{b,L}
    =
    \frac{bR(x)-R(bx)}{\ln b}.
    \label{eq:tail-R-identity}
\end{equation}
Using
\eqref{eq:R-stirling},
\begin{align}
    bR(x)-R(bx)
    &=
    \frac{b-b^{-1}}{12x}
    -
    \frac{b-b^{-3}}{360x^3}
    \notag\\
    &\quad+
    b\rho(x)-\rho(bx).
    \label{eq:bR-minus-R}
\end{align}

Since
\[
    x=b^{L-1},
\]
we have
\begin{equation}
    \frac{b-b^{-1}}{x}
    =
    (b^2-1)b^{-L}
    \label{eq:first-tail-conversion}
\end{equation}
and
\begin{equation}
    \frac{b-b^{-3}}{x^3}
    =
    (b^4-1)b^{-3L}.
    \label{eq:second-tail-conversion}
\end{equation}
Thus
\eqref{eq:tail-R-identity}
becomes
\eqref{eq:intro-tail-expansion}, with
\begin{equation}
    \varepsilon_{b,L}
    =
    \frac{
        b\rho(x)-\rho(bx)
    }{\ln b}.
    \label{eq:epsilon-definition}
\end{equation}

Because \(b\ge2\), one has \(\ln b>0\).
Using
\eqref{eq:stirling-rho-bound} and the triangle inequality,
\begin{align}
    |\varepsilon_{b,L}|
    &\le
    \frac{
        b\rho(x)+\rho(bx)
    }{\ln b}
    \notag\\
    &\le
    \frac{1}{1260\ln b}
    \left(
        \frac{b}{x^5}
        +
        \frac{1}{b^5x^5}
    \right).
    \label{eq:epsilon-bound-intermediate}
\end{align}
Finally,
\[
    x=b^{L-1}
\]
implies
\[
    \frac{b}{x^5}
    =
    b^6b^{-5L},
    \qquad
    \frac{1}{b^5x^5}
    =
    b^{-5L},
\]
and hence
\[
    |\varepsilon_{b,L}|
    \le
    \frac{b^6+1}{1260\ln b}\,b^{-5L}.
\]
\end{proof}

\begin{remark}
The absence of \(b^{-2L}\) and \(b^{-4L}\) terms in
\eqref{eq:intro-tail-expansion} reflects the odd inverse powers
in the Stirling series for \(\ln\Gamma(x)\).
\end{remark}


\subsection{Finite-prefix average}
\label{subsec:finite-prefix-average-refined}

The preceding theorem sharpens
\eqref{eq:model-average-leading} immediately.

\begin{corollary}
\label{cor:finite-prefix-average-refined}
For every \(b\ge2\) and \(L\ge1\),
\begin{equation}
\boxed{
\begin{aligned}
    \mathbb E[A_L^{(b)}]
    &=
    \frac{b-1}{2}
    -
    \frac{c_b}{L}
    \\
    &\quad+
    \frac{b^2-1}{12\ln b}
    \frac{b^{-L}}{L}
    -
    \frac{b^4-1}{360\ln b}
    \frac{b^{-3L}}{L}
    +
    \frac{\varepsilon_{b,L}}{L},
\end{aligned}
}
\label{eq:finite-prefix-average-refined}
\end{equation}
where
\[
    |\varepsilon_{b,L}|
    \le
    \frac{b^6+1}{1260\ln b}\,b^{-5L}.
\]
\end{corollary}

\begin{proof}
From
\eqref{eq:model-average-exact},
\[
    \mathbb E[A_L^{(b)}]
    =
    \frac{b-1}{2}
    -
    \frac{c_{b,L}}{L}.
\]
Write
\[
    c_{b,L}
    =
    c_b-(c_b-c_{b,L})
\]
and apply
Theorem~\ref{thm:intro-tail}.
\end{proof}

The expansion
\eqref{eq:finite-prefix-average-refined}
shows that the \(L^{-1}\) correction determined by \(c_b\) is separated
from exponentially small finite-prefix effects.  This distinction will
be useful in the deterministic setting: the next section asks whether
the same \(O(1)\) cumulative correction survives when the logarithmic
mantissa integral is replaced by an orbit average along
\(\{n\log_b a\}\).

% ============================================================
% Section 5
% ============================================================

\section{Deterministic exponential sequences and discrepancy transfer}
\label{sec:deterministic-transfer}

We now pass from the logarithmic mantissa probability model to actual
deterministic exponential sequences.

Let \(a,b\ge2\) be integers and assume that they are multiplicatively
independent.  Equivalently,
\[
    \alpha=\log_b a
\]
is irrational.  Define
\begin{equation}
    \theta_n=\{n\alpha\},
    \qquad n\ge1.
    \label{eq:theta-n-deterministic}
\end{equation}
For \(\theta\in[0,1)\), we write
\(D_j^{(b)}(b^\theta)\) for the \(j\)-th significant
base-\(b\) digit of \(b^\theta\).
Writing
\[
    n\alpha=\lfloor n\alpha\rfloor+\theta_n,
\]
we have
\begin{equation}
    a^n
    =
    b^{\lfloor n\alpha\rfloor}
    b^{\theta_n}.
    \label{eq:power-mantissa-decomposition}
\end{equation}
Thus the significant digits of \(a^n\) are determined by the orbit
\((\theta_n)\).

The number of base-\(b\) digits of \(a^n\) is
\begin{equation}
    L_n
    =
    \lfloor n\alpha\rfloor+1.
    \label{eq:L-n-deterministic}
\end{equation}

Throughout this section, \(D_N(\alpha)\) denotes the ordinary interval
discrepancy of the first \(N\) points of the rotation:
\begin{equation}
\boxed{
    D_N(\alpha)
    =
    \sup_{0\le u<v\le1}
    \left|
        \frac1N
        \#\{1\le n\le N:\theta_n\in[u,v)\}
        -(v-u)
    \right|.
}
\label{eq:ordinary-discrepancy}
\end{equation}
We use ordinary discrepancy rather than star discrepancy because the
digit events below are naturally unions of arbitrary half-open
subintervals of \([0,1)\).


\subsection{Significant-digit cylinders}
\label{subsec:digit-cylinders}

We first describe the interval geometry of a fixed significant-digit
position.

For \(j\ge2\), put
\[
    M=b^{j-2}.
\]
Fix
\[
    r\in\{0,\ldots,b-1\}.
\]
The condition that the \(j\)-th significant digit of \(b^\theta\) is
\(r\) is equivalent to the \(j\)-digit integer prefix having the form
\[
    bq+r,
    \qquad
    M\le q\le bM-1.
\]
For each such \(q\), define
\begin{equation}
    I_{q,r}^{(j)}
    =
    \left[
        \log_b
        \left(
            \frac{bq+r}{b^{j-1}}
        \right),
        \log_b
        \left(
            \frac{bq+r+1}{b^{j-1}}
        \right)
    \right).
    \label{eq:digit-cylinder}
\end{equation}

\begin{lemma}[Digit-cylinder decomposition]
\label{lem:digit-cylinder-decomposition}
For \(j\ge2\) and
\(r\in\{0,\ldots,b-1\}\), the event
\[
    D_j^{(b)}(b^\theta)=r
\]
is the disjoint union
\begin{equation}
    \bigcup_{q=b^{j-2}}^{b^{j-1}-1}
    I_{q,r}^{(j)}.
    \label{eq:digit-cylinder-union}
\end{equation}
In particular, it consists of exactly
\begin{equation}
\boxed{
    K_j=(b-1)b^{j-2}
}
\label{eq:K-j}
\end{equation}
half-open intervals.
\end{lemma}

\begin{proof}
The first \(j\) significant base-\(b\) digits of \(b^\theta\) form
the integer
\[
    \left\lfloor b^{j-1+\theta}\right\rfloor.
\]
Its final digit is \(r\) if and only if it equals \(bq+r\) for some
\[
    b^{j-2}\le q\le b^{j-1}-1.
\]
For a fixed \(q\),
\[
    \left\lfloor b^{j-1+\theta}\right\rfloor=bq+r
\]
is equivalent to
\[
    bq+r
    \le
    b^{j-1+\theta}
    <
    bq+r+1,
\]
which is exactly
\[
    \theta\in I_{q,r}^{(j)}.
\]
Distinct prefix integers give disjoint half-open intervals.  The
number of admissible \(q\) is
\[
    b^{j-1}-b^{j-2}
    =
    (b-1)b^{j-2}.
\]
\end{proof}

The total length of the union in
\eqref{eq:digit-cylinder-union} is
\[
    \mathbb P(D_j^{(b)}=r),
\]
the logarithmic-mantissa probability from
Section~\ref{subsec:higher-digit-laws}.

For the leading digit, the corresponding cylinders are simply
\begin{equation}
    I_d^{(1)}
    =
    [\log_b d,\log_b(d+1)),
    \qquad
    d=1,\ldots,b-1.
    \label{eq:first-digit-cylinder}
\end{equation}


\subsection{Fixed-position discrepancy transfer}
\label{subsec:fixed-position-transfer}

Lemma~\ref{lem:digit-cylinder-decomposition} gives an immediate
quantitative comparison between the deterministic orbit and the
logarithmic mantissa model.

\begin{proposition}[Fixed-position transfer]
\label{prop:fixed-position-transfer}
For \(j\ge2\) and
\(r\in\{0,\ldots,b-1\}\),
\begin{equation}
\boxed{
\left|
    \frac1N
    \#\{1\le n\le N:
        D_j^{(b)}(b^{\theta_n})=r\}
    -
    \mathbb P(D_j^{(b)}=r)
\right|
\le
    (b-1)b^{j-2}D_N(\alpha).
}
\label{eq:fixed-position-transfer}
\end{equation}
\end{proposition}

\begin{proof}
The digit event is the disjoint union of
\(K_j=(b-1)b^{j-2}\) intervals from
Lemma~\ref{lem:digit-cylinder-decomposition}.
For each interval \(I\), the definition of ordinary discrepancy gives
\[
    \left|
        \frac1N
        \#\{1\le n\le N:\theta_n\in I\}
        -
        |I|
    \right|
    \le
    D_N(\alpha).
\]
Summing over the \(K_j\) disjoint components gives
\eqref{eq:fixed-position-transfer}.
\end{proof}

To translate this into a statement about the actual integer \(a^n\),
we must account for the finitely many early exponents for which the
\(j\)-th significant digit does not yet exist.

Define
\begin{equation}
    e_j(N)
    =
    \#\{1\le n\le N:L_n<j\}.
    \label{eq:e-j-definition}
\end{equation}
From
\[
    L_n<j
    \iff
    n\alpha<j-1,
\]
one has the exact formula
\begin{equation}
    e_j(N)
    =
    \min
    \left\{
        N,\,
        \max
        \left(
            0,\,
            \left\lceil
                \frac{j-1}{\alpha}
            \right\rceil-1
        \right)
    \right\}.
    \label{eq:e-j-exact}
\end{equation}
In particular,
\begin{equation}
    e_j(N)=O_{a,b}(j).
    \label{eq:e-j-big-O}
\end{equation}

\begin{corollary}
\label{cor:actual-fixed-position-transfer}
For \(j\ge2\),
\begin{equation}
\begin{aligned}
&\left|
    \frac1N
    \#\{
        1\le n\le N:
        L_n\ge j,\,
        D_j^{(b)}(a^n)=r
    \}
    -
    \mathbb P(D_j^{(b)}=r)
\right|
\\
&\hspace{35mm}
\le
    (b-1)b^{j-2}D_N(\alpha)
    +
    \frac{e_j(N)}{N}.
\end{aligned}
\label{eq:actual-fixed-position-transfer}
\end{equation}
\end{corollary}

\begin{proof}
Whenever \(L_n\ge j\), the first \(j\) significant digits of
\(a^n\) are the first \(j\) significant digits determined by
\(b^{\theta_n}\).  Starting from
Proposition~\ref{prop:fixed-position-transfer}, deleting the
\(e_j(N)\) short powers changes the normalized count by at most
\(e_j(N)/N\).
\end{proof}

Combining
Corollary~\ref{cor:actual-fixed-position-transfer} with
Proposition~\ref{prop:deep-position-expansion} gives, for fixed \(b\),
\begin{equation}
\begin{aligned}
    &\frac1N
    \#\{
        1\le n\le N:
        L_n\ge j,\,
        D_j^{(b)}(a^n)=r
    \}
    \\
    &\qquad=
    \frac1b
    +
    \frac{(b-1)(b-1-2r)}
         {2b\ln b}\,b^{-j}
    +
O_b\!\left(
    b^{-2j}
    +
    b^jD_N(\alpha)
\right)
+
O_{a,b}\!\left(\frac{j}{N}\right).
    \label{eq:fixed-position-combined}
\end{aligned}
\end{equation}
This estimate is useful for orientation, but the main result below
will sum the first \(J\) positions simultaneously.


\subsection{The leading-window observable}
\label{subsec:leading-window-observable}

For \(J\ge1\), define
\begin{equation}
    G_J(\theta)
    =
    s_b
    \left(
        \left\lfloor
            b^{J-1+\theta}
        \right\rfloor
    \right),
    \qquad
    0\le\theta<1.
    \label{eq:G-J-definition}
\end{equation}
This is the sum of the first \(J\) significant digits of \(b^\theta\).

\begin{lemma}[Exact leading-window coding]
\label{lem:leading-window-coding}
If \(L_n\ge J\), then
\begin{equation}
\boxed{
    G_J(\theta_n)
    =
    \sum_{j=1}^{J}
    D_j^{(b)}(a^n).
}
\label{eq:leading-window-coding}
\end{equation}
\end{lemma}

\begin{proof}
From
\eqref{eq:power-mantissa-decomposition},
\[
    a^n
    =
    b^{L_n-1}b^{\theta_n}.
\]
Multiplying the normalized significand \(b^{\theta_n}\) by
\(b^{J-1}\) and taking the integer part extracts exactly the first
\(J\) significant digits, provided those digits exist.  Thus
\[
    \left\lfloor
        b^{J-1+\theta_n}
    \right\rfloor
\]
is the \(J\)-digit integer prefix of \(a^n\), and its digit sum is
the right-hand side of
\eqref{eq:leading-window-coding}.
\end{proof}

The logarithmic-mantissa average of \(G_J\) is already known exactly
from Section~\ref{sec:finite-bias}.

\begin{lemma}[Integral of the leading-window observable]
\label{lem:G-J-integral}
For every \(J\ge1\),
\begin{equation}
\boxed{
    \int_0^1 G_J(\theta)\,d\theta
    =
    J\frac{b-1}{2}
    -
    c_{b,J}.
}
\label{eq:G-J-integral}
\end{equation}
\end{lemma}

\begin{proof}
The random integer
\[
    \left\lfloor b^{J-1+\Theta}\right\rfloor
\]
is precisely \(M_J\) from
Section~\ref{subsec:finite-significands}.  Hence
\[
    \int_0^1 G_J(\theta)\,d\theta
    =
    \mathbb E[s_b(M_J)].
\]
By the definition of \(c_{b,J}\),
\[
    \mathbb E[s_b(M_J)]
    =
    J\frac{b-1}{2}-c_{b,J}.
\]
\end{proof}


\subsection{A finite discrepancy bound for a leading window}
\label{subsec:finite-leading-window-bound}

Define the deterministic leading-window average
\begin{equation}
    \mathcal S_{N,J}(a;b)
    =
    \frac1N
    \sum_{\substack{1\le n\le N\\L_n\ge J}}
    \sum_{j=1}^{J}
    D_j^{(b)}(a^n).
    \label{eq:S-NJ-main}
\end{equation}

Before deleting the short powers, consider
\[
    T_{N,J}
    =
    \frac1N
    \sum_{n=1}^{N}G_J(\theta_n).
\]
Since
\[
    G_J(\theta)
    =
    \sum_{j=1}^{J}D_j^{(b)}(b^\theta),
\]
we may estimate each digit position by its interval cylinders.

For the first digit,
\eqref{eq:first-digit-cylinder} and ordinary discrepancy give
\begin{equation}
\begin{aligned}
    \left|
        \frac1N\sum_{n=1}^{N}
        D_1^{(b)}(b^{\theta_n})
        -
        \mathbb E[D_1^{(b)}]
    \right|
    &\le
    \sum_{d=1}^{b-1}d\,D_N(\alpha)
    \\
    &=
    \frac{b(b-1)}{2}D_N(\alpha).
\end{aligned}
\label{eq:first-digit-weighted-discrepancy}
\end{equation}

For \(j\ge2\), there are
\[
    K_j=(b-1)b^{j-2}
\]
components for each digit \(r\).  Hence
\begin{align}
    &\left|
        \frac1N\sum_{n=1}^{N}
        D_j^{(b)}(b^{\theta_n})
        -
        \mathbb E[D_j^{(b)}]
    \right|
    \notag\\
    &\qquad\le
    K_jD_N(\alpha)
    \sum_{r=0}^{b-1}r
    \notag\\
    &\qquad=
    \frac{b(b-1)}{2}
    (b-1)b^{j-2}
    D_N(\alpha).
    \label{eq:higher-digit-weighted-discrepancy}
\end{align}

Summing
\eqref{eq:first-digit-weighted-discrepancy}
and
\eqref{eq:higher-digit-weighted-discrepancy}
over \(1\le j\le J\), the geometric series collapses:
\begin{align}
    &\frac{b(b-1)}2
    +
    \sum_{j=2}^{J}
    \frac{b(b-1)}2
    (b-1)b^{j-2}
    \notag\\
    &\qquad=
    \frac{b(b-1)}2
    \left(
        1+
        (b-1)\sum_{j=2}^{J}b^{j-2}
    \right)
    \notag\\
    &\qquad=
    \frac{b-1}{2}b^J.
    \label{eq:weighted-component-total}
\end{align}
Therefore
\begin{equation}
    \left|
        T_{N,J}
        -
        \int_0^1G_J(\theta)\,d\theta
    \right|
    \le
    \frac{b-1}{2}
    b^J D_N(\alpha).
    \label{eq:T-NJ-discrepancy}
\end{equation}

We can now prove the finite deterministic transfer theorem announced
in the introduction.

\begin{proof}[Proof of
Theorem~\ref{thm:intro-deterministic-finite}]
By
Lemma~\ref{lem:leading-window-coding},
\[
    \mathcal S_{N,J}(a;b)
\]
agrees with \(T_{N,J}\) on every exponent satisfying \(L_n\ge J\).
The two quantities differ only at the
\[
    e_J(N)
\]
short powers.

Since
\[
    0\le G_J(\theta)\le J(b-1),
\]
we have
\begin{equation}
    \left|
        \mathcal S_{N,J}(a;b)-T_{N,J}
    \right|
    \le
    \frac{(b-1)J\,e_J(N)}{N}.
    \label{eq:short-power-cost}
\end{equation}
Combining
\eqref{eq:T-NJ-discrepancy},
\eqref{eq:short-power-cost}, and
Lemma~\ref{lem:G-J-integral} gives
\[
\boxed{
\begin{aligned}
    \left|
        \mathcal S_{N,J}(a;b)
        -
        \left(
            J\frac{b-1}{2}-c_{b,J}
        \right)
    \right|
    \le{}&
    \frac{b-1}{2}
    b^JD_N(\alpha)
    \\
    &+
    \frac{(b-1)J\,e_J(N)}{N}.
\end{aligned}
}
\]
\end{proof}


\subsection{The second-order deterministic term}
\label{subsec:second-order-deterministic}

The finite bound immediately isolates the cumulative Benford bias as
the bounded correction term.

\begin{proof}[Proof of
Theorem~\ref{thm:intro-second-order}]
From
Theorem~\ref{thm:intro-deterministic-finite},
\[
\begin{aligned}
    \mathcal S_{N,J}(a;b)
    =
    J\frac{b-1}{2}
    -
    c_{b,J}
    +
    O_b\!\left(b^JD_N(\alpha)\right)
    +
    O_{a,b}\!\left(\frac{J^2}{N}\right),
\end{aligned}
\label{eq:pre-D6}
\]
where we used
\[
    e_J(N)=O_{a,b}(J).
\]

By
Theorem~\ref{thm:intro-tail},
\[
    c_{b,J}
    =
    c_b+O_b(b^{-J}).
\]
Consequently,
\begin{equation}
\boxed{
\begin{aligned}
    \mathcal S_{N,J}(a;b)
    =
    J\frac{b-1}{2}
    -
    c_b
    &+
    O_b(b^{-J})
    \\
    &+
    O_b\!\left(b^JD_N(\alpha)\right)
    +
    O_{a,b}\!\left(\frac{J^2}{N}\right).
\end{aligned}
}
\label{eq:D6-main}
\end{equation}

If
\[
    J=J(N)\longrightarrow\infty,
\]
\[
    b^{J(N)}D_N(\alpha)\longrightarrow0,
\]
and
\[
    \frac{J(N)^2}{N}\longrightarrow0,
\]
then every error term in
\eqref{eq:D6-main} tends to zero.  Hence
\[
\boxed{
    \mathcal S_{N,J(N)}(a;b)
    =
    J(N)\frac{b-1}{2}
    -
    c_b
    +
    o(1).
}
\]
\end{proof}

Equivalently,
\begin{equation}
\boxed{
    J(N)\frac{b-1}{2}
    -
    \mathcal S_{N,J(N)}(a;b)
    \longrightarrow
    c_b.
}
\label{eq:cb-deterministic-limit}
\end{equation}
Thus the same constant obtained by integrating against the logarithmic
mantissa measure survives as the bounded correction in a deterministic
orbit average.


\subsection{A power-saving discrepancy input}
\label{subsec:power-saving-discrepancy}

The preceding theorem is formulated conditionally in terms of
\(D_N(\alpha)\).  For a fixed multiplicatively independent pair
\((a,b)\), classical Diophantine-discrepancy theory supplies the
required quantitative input.

Fuchs, building on a result of Kiss and Tichy and Baker's theory of
linear forms in logarithms, records the following power-saving
discrepancy estimate for irrational quotients of logarithms
\cite[Lemma~2]{Fuchs2002}.

\begin{proposition}[Classical power-saving discrepancy]
\label{prop:power-saving-discrepancy}
Let \(a,b\ge2\) be fixed multiplicatively independent integers.  Then
there exists a constant
\[
    \gamma_{a,b}>0
\]
such that
\begin{equation}
\boxed{
    D_N(\log_b a)
    \ll_{a,b}
    N^{-\gamma_{a,b}}.
}
\label{eq:power-saving-discrepancy}
\end{equation}
\end{proposition}

\begin{remark}[External input]
\label{rem:discrepancy-external-input}
Proposition~\ref{prop:power-saving-discrepancy} is not a result of this
paper.  Fuchs states the estimate for
\[
    \alpha=\log_A Q
\]
and the sequence \(-n\alpha\), as a special case of a more general
result of Kiss and Tichy, and notes that the proof follows from the
Erd\H{o}s--Tur\'an inequality together with Baker's theorem on linear
forms in logarithms
\cite{Fuchs2002,KissTichy1989I,KissTichy1989II}.
Taking
\[
    A=b,
    \qquad
    Q=a
\]
gives \(\alpha=\log_b a\).
The sign is immaterial for ordinary discrepancy, since reflection
modulo one maps intervals to intervals of the same length.
\end{remark}


\subsection{Logarithmically growing leading windows}
\label{subsec:logarithmic-leading-windows}

We can now remove the discrepancy hypothesis from the main
second-order statement for each fixed pair \((a,b)\).

\begin{proof}[Proof of
Corollary~\ref{cor:intro-log-window}]
Fix
\[
    0<\kappa<\gamma_{a,b}
\]
and define
\[
    J_N
    =
    \left\lfloor
        \kappa\log_bN
    \right\rfloor.
\]
Then
\[
    J_N\longrightarrow\infty
\]
and
\[
    \frac{J_N^2}{N}
    =
    O\!\left(
        \frac{(\log N)^2}{N}
    \right)
    \longrightarrow0.
\]
Moreover,
\[
    b^{J_N}
    \le
    N^\kappa.
\]
Using
Proposition~\ref{prop:power-saving-discrepancy},
\[
    b^{J_N}D_N(\log_ba)
    \ll_{a,b}
    N^{\kappa-\gamma_{a,b}}
    \longrightarrow0.
\]
The hypotheses of
Theorem~\ref{thm:intro-second-order}
are therefore satisfied, and
\[
\boxed{
    \mathcal S_{N,J_N}(a;b)
    =
    J_N\frac{b-1}{2}
    -
    c_b
    +
    o(1).
}
\]
\end{proof}

The proof also gives an explicit qualitative rate.

\begin{corollary}[Quantitative logarithmic-window form]
\label{cor:quantitative-log-window}
Under the hypotheses above,
\begin{equation}
\begin{aligned}
    \mathcal S_{N,J_N}(a;b)
    =
    J_N\frac{b-1}{2}
    -
    c_b
    +
    O_{a,b,\kappa}
    \left(
        N^{-\kappa}
        +
        N^{-(\gamma_{a,b}-\kappa)}
        +
        \frac{(\log N)^2}{N}
    \right).
    \label{eq:quantitative-log-window}
\end{aligned}
\end{equation}
\end{corollary}

\begin{proof}
Since
\[
    J_N
    \ge
    \kappa\log_bN-1,
\]
we have
\[
    b^{-J_N}
    \ll_b N^{-\kappa}.
\]
The other two terms follow directly from
\eqref{eq:D6-main} and
\eqref{eq:power-saving-discrepancy}.
\end{proof}

For example, choosing
\[
    \kappa=\frac{\gamma_{a,b}}2
\]
balances the first two power-saving errors and yields
\begin{equation}
    \mathcal S_{N,J_N}(a;b)
    =
    J_N\frac{b-1}{2}
    -
    c_b
    +
    O_{a,b}
    \left(
        N^{-\gamma_{a,b}/2}
        +
        \frac{(\log N)^2}{N}
    \right),
    \label{eq:balanced-log-window}
\end{equation}
with
\[
    J_N
    =
    \left\lfloor
        \frac{\gamma_{a,b}}2
        \log_bN
    \right\rfloor.
\]


\subsection{Relation to earlier growing-window results}
\label{subsec:fuchs-comparison}

The logarithmic growth of the controlled window is not itself new.
Fuchs studied digital blocks in growing regions of exponential
sequences and used discrepancy and exponential-sum estimates to obtain
first-order asymptotic digit statistics
\cite{Fuchs2002}.
The discrepancy machinery in
Proposition~\ref{prop:power-saving-discrepancy}
likewise belongs to that earlier literature.

The point of
\eqref{eq:cb-deterministic-limit}
is instead second-order.
The first-order term
\[
    J\frac{b-1}{2}
\]
reflects the asymptotically uniform bulk digit distribution.
After this term is removed, the remaining bounded contribution is not
discarded inside an \(O(1)\) error: it converges to
\[
    -c_b.
\]
Thus
\[
    c_b
    =
    \frac b2
    -
    \frac{b-1}{2}\log_b(2\pi)
\]
acts as a boundary correction generated by the finitely biased
significant positions near the leading edge.

To the best of our knowledge, this explicit constant term has not
previously been recorded in the growing-window theory for exponential
sequences.


\subsection{Scope of the deterministic result}
\label{subsec:deterministic-scope}

It is important to distinguish the leading-window theorem from a
statement about the complete base-\(b\) digit sum of \(a^n\).

For \(n\) of order \(N\),
\[
    L_n
    =
    n\log_ba+O(1)
    \asymp_{a,b}N,
\]
whereas the unconditional window obtained above has only
\[
    J_N=O_{a,b}(\log N)
\]
positions.
Hence
\[
    \frac{J_N}{L_n}\longrightarrow0
\]
for typical \(n\) in the averaging range.

Accordingly, the present argument controls a vanishing fraction of the
complete digit expansion.  It does not establish a mean law, variance
law, central limit theorem, mixing property, or independence statement
for the full digit sum \(s_b(a^n)\).  Extending absolute-error control
from a logarithmic leading window to a positive proportion of the
digits would require substantially stronger input than the discrepancy
estimate used here.

% ============================================================
% Section 6
% ============================================================

\section{Computational verification and reproducibility}
\label{sec:verification}

The results of Sections~\ref{sec:finite-bias}--\ref{sec:deterministic-transfer}
are proved analytically and do not depend on numerical computation.
Nevertheless, because several analytic identities involve substantial
cancellation and the deterministic estimates require delicate
indexing, boundary handling, and discrepancy bookkeeping, we carried
out a separate adversarial computational audit of the full theorem
package.

The purpose of this audit was not to supply numerical evidence in place
of proof.  Rather, it was designed to detect indexing errors,
normalization mistakes, hidden boundary cases, numerical instability,
or disagreements between mathematically equivalent formulations.

\subsection{Audit structure}
\label{subsec:audit-structure}

The verification package treats the principal statements as sixteen
separate claims, denoted internally by
\[
    \mathrm{A1},\ldots,\mathrm{A10},
    \qquad
    \mathrm{D1},\ldots,\mathrm{D6}.
\]
The analytic claims include the exact significant-digit formulas, the
finite Benford-weighted digit-sum identity, the closed expression for
\(c_b\), the deep-position expansions, and the finite-length Stirling
remainder estimate.  The deterministic claims include the digit-cylinder
decomposition, leading-window coding, and the finite discrepancy
transfer estimates.

Where possible, the same mathematical quantity was recomputed through
independent representations.  In particular, the audit compared:

\begin{itemize}
    \item direct logarithmic interval sums with Gamma-function formulas;

    \item position-wise digit expectations with direct enumeration of
    finite integer significands;

    \item the finite cumulative bias obtained from individual digits
    with the factorial identity of
    Theorem~\ref{thm:intro-finite-digit-sum};

    \item mantissa-based digit extraction with exact integer base-\(b\)
    digit extraction for powers \(a^n\);

    \item the theoretical discrepancy bounds with explicit orbit
    counts for finite test cases.
\end{itemize}

All comparisons involving Gamma functions or deep Stirling tails were
performed at arbitrary precision rather than ordinary machine
floating-point precision.


\subsection{Verification results}
\label{subsec:verification-results}

No candidate identity in the final theorem package was falsified.
All sixteen claims passed the final audit.

For the analytic part, the finite prefix probabilities, Gamma
identities, factorial identity, and closed formula for \(c_b\) agreed
across independent computational routes to high precision.  The
deep-position coefficients in
Proposition~\ref{prop:deep-position-expansion} and
Corollary~\ref{cor:mean-digit-expansion} were also checked by scaled
residual experiments across several bases.

The finite-tail theorem
Theorem~\ref{thm:intro-tail} received a separate
high-precision stress test.  Because the relevant expression involves
subtraction of large log-Gamma terms, low working precision can produce
spurious violations of the remainder estimate.  The final experiments
therefore used 220 decimal digits of precision.  The numerical tests
were consistent with the rigorous bound
\[
    |\varepsilon_{b,L}|
    \le
    \frac{b^6+1}{1260\ln b}\,b^{-5L}.
\]
These experiments are supplementary only: the proof of the bound is
the analytic positive-real Stirling remainder argument in
Section~\ref{subsec:cumulative-tail-expansion}.

For the deterministic part, exact integer digit extraction was compared
against the mantissa representation for thousands of powers and their
tested leading positions, with no disagreement after explicit guarding
of exact and near-exact integer boundaries.  The coefficient
\[
    \frac{b-1}{2}b^J
\]
in the finite leading-window discrepancy bound was also independently
reconstructed from the weighted interval-component counts and tested
against direct finite orbit computations.

The final automated test suite contained sixteen tests, all of which
passed.


\subsection{Numerical boundary issues}
\label{subsec:numerical-boundaries}

Two numerical issues encountered during the audit are worth recording.

First, evaluating
\[
    \left\lfloor b^{J-1+\theta}\right\rfloor
\]
directly with finite-precision floating-point arithmetic can be
unreliable when the exact value is an integer or lies extremely close
to one.  Exact integer arithmetic was therefore treated as the ground
truth for digit extraction, and high-precision nearest-integer checks
were used at mantissa boundaries.

Second, direct evaluation of the finite-bias formula
\[
    \ln\Gamma(b^L)-b\ln\Gamma(b^{L-1})
\]
can suffer severe cancellation for large \(L\).  High-precision
arithmetic was used in all tests involving the exponentially small
tail \(c_b-c_{b,L}\).

These observations concern implementation only and do not alter the
analytic statements proved above.


\subsection{Reproducibility archive}
\label{subsec:reproducibility-archive}

The computational audit is organized as a separate reproducibility
repository containing:

\begin{itemize}
    \item an authoritative theorem specification matching the claims
    proved in the manuscript;

    \item independent implementations of the principal formulas;

    \item automated unit and adversarial tests;

    \item high-precision residual experiments for the asymptotic
    statements;

    \item deterministic digit-extraction and discrepancy experiments;

    \item machine-readable claim-status and numerical-residual tables;

    \item instructions sufficient to recreate the final verification
    run.
\end{itemize}

The verification code is intentionally separated from the proofs in
this paper.  Its role is to make the formulas independently checkable
and to document numerical edge cases, while the mathematical validity
of the main theorems rests entirely on the arguments in
Sections~\ref{sec:finite-bias}--\ref{sec:deterministic-transfer}.

% TODO before public submission:
% Insert the permanent GitHub/Zenodo archival citation here once the
% verification repository has been finalized and deposited.

% ============================================================
% Section 7
% ============================================================

\section{Discussion, limitations, and open problems}
\label{sec:discussion}

The results of this paper identify a precise second-order contribution
to the mean sum of a growing block of leading significant digits in an
exponential sequence.  The main mechanism can be summarized as follows:
the early significant positions are not individually uniform under the
logarithmic mantissa measure, their cumulative signed deviation has the
finite limit
\[
    c_b
    =
    \frac b2
    -
    \frac{b-1}{2}\log_b(2\pi),
\]
and quantitative equidistribution transfers this same bounded correction
to a deterministic leading-window average.

This section records the limits of that argument and several natural
directions in which the present results might be extended.


\subsection{Leading windows versus complete digit sums}
\label{subsec:leading-vs-complete}

The principal deterministic result concerns
\[
    \mathcal S_{N,J}(a;b)
    =
    \frac1N
    \sum_{\substack{1\le n\le N\\L_n\ge J}}
    \sum_{j=1}^{J}D_j^{(b)}(a^n),
\]
with
\[
    J=J(N)
\]
chosen so that
\[
    b^J D_N(\log_b a)\longrightarrow0.
\]
For the unconditional fixed-pair corollary obtained from a power-saving
discrepancy estimate, this gives
\[
    J_N=O_{a,b}(\log N).
\]

By contrast, for \(n\asymp N\), the complete number of base-\(b\)
digits of \(a^n\) is
\[
    L_n
    =
    n\log_b a+O(1)
    \asymp_{a,b}N.
\]
Hence the controlled leading block satisfies
\[
    \frac{J_N}{L_n}\longrightarrow0.
\]

The theorem therefore does not determine the asymptotic behavior of
the complete sum of digits
\[
    s_b(a^n)
    =
    \sum_{j=1}^{L_n}D_j^{(b)}(a^n),
\]
nor of the normalized complete-digit average
\[
    \frac{s_b(a^n)}{L_n}.
\]
Controlling a logarithmic number of significant positions is sufficient
to retain the finite boundary constant \(c_b\), but it controls only a
vanishing fraction of the full expansion.

This distinction is essential: the second-order theorem proved here
should not be interpreted as a complete digit-sum theorem for
exponential sequences.


\subsection{The barrier imposed by discrepancy}
\label{subsec:discrepancy-barrier}

The finite deterministic estimate
\[
    \mathcal S_{N,J}(a;b)
    =
    J\frac{b-1}{2}
    -
    c_b
    +
    O_b(b^{-J})
    +
    O_b\!\left(b^JD_N(\log_b a)\right)
    +
    O_{a,b}\!\left(\frac{J^2}{N}\right)
\]
makes the present limitation explicit.

The complexity of the first \(J\) significant positions grows
essentially like \(b^J\).  Ordinary discrepancy must therefore decrease
fast enough to compensate for this exponential growth.  A polynomial
bound
\[
    D_N(\log_b a)\ll N^{-\gamma}
\]
naturally leads to a logarithmic admissible range
\[
    J\lesssim\gamma\log_bN.
\]

Thus one obvious open problem is to determine whether stronger
arithmetic or harmonic analysis can control significantly larger
leading windows while retaining an absolute \(o(1)\) error after
subtracting the first-order term.

The difficulty is stronger than proving first-order digit
equidistribution.  A statement of the form
\[
    \mathcal S_{N,J}
    =
    J\frac{b-1}{2}
    +
    o(J)
\]
does not determine a bounded second-order term.  To identify
\(-c_b\), the total error must instead be controlled on the absolute
scale \(o(1)\).


\subsection{Beyond logarithmic windows}
\label{subsec:beyond-logarithmic-windows}

A natural first extension would be to replace
\[
    J_N\asymp\log N
\]
by a faster-growing sequence while preserving
\[
    \mathcal S_{N,J_N}(a;b)
    =
    J_N\frac{b-1}{2}
    -
    c_b
    +
    o(1).
\]

There are at least two conceptually different ways such an improvement
could arise.  One could seek sharper discrepancy estimates tailored to
the specific rotation
\[
    \{n\log_b a\},
\]
or one could exploit more structure in the digit observable itself,
rather than bounding all of its interval components independently.

The present proof deliberately uses ordinary discrepancy in a robust
but coarse way.  It does not exploit possible cancellation between the
many digit cylinders.  Whether such cancellation can enlarge the
second-order window remains open.


\subsection{The complete average-digit problem}
\label{subsec:complete-average-digit}

A substantially harder problem is to understand
\begin{equation}
    A_{a,b}(n)
    =
    \frac{s_b(a^n)}{L_n}.
    \label{eq:complete-average-digit}
\end{equation}
The logarithmic mantissa model suggests the uniform digit mean
\[
    \frac{b-1}{2}
\]
as the natural center, but the present results do not prove
\[
    A_{a,b}(n)\to\frac{b-1}{2}
\]
pointwise, in mean over \(n\), or in any stronger sense.

The obstacle is structural.  For each \(n\), the statistic
\eqref{eq:complete-average-digit} depends on a number of digit
positions that grows linearly with \(n\).  Standard equidistribution of
\[
    \{n\log_b a\}
\]
controls every fixed observable, while the complete digit average
corresponds to a family of observables whose complexity increases with
\(n\).

Bridging this varying-observable gap appears to require ideas beyond
the fixed-position and logarithmic-window methods used here.


\subsection{Fluctuations and dependence}
\label{subsec:fluctuations-open}

The present paper concerns expectations, cumulative bias, and
deterministic orbit averages.  It does not establish a variance law or
a limiting distribution for digit sums of \(a^n\).

In particular, no assumption of independence among significant digits
is made.  Irrational rotations are equidistributed but are not mixing,
and fixed-position marginal equidistribution does not imply temporal
independence, Gaussian fluctuations, or a central limit theorem for a
changing digit observable.

Natural future questions include whether appropriately normalized
leading-window or complete-digit statistics possess limiting variance
laws, whether arithmetic properties of \(a\) and \(b\) influence those
fluctuations, and whether any probabilistic limit theorem is available
under additional hypotheses.

These questions are deliberately left separate from the deterministic
second-order theory proved above.


\subsection{Other digital functions}
\label{subsec:other-digital-functions}

The finite identity in
Theorem~\ref{thm:intro-finite-digit-sum} arises from two ingredients:
the logarithmic Benford weight
\[
    \log_b\!\left(\frac{m+1}{m}\right)
\]
and the radix recursion
\[
    s_b(bq+r)=s_b(q)+r.
\]
This suggests a broader question.

Let \(f\) be a digital or \(b\)-additive function.  One may ask whether
finite logarithmically weighted expectations
\[
    \sum_{m=b^{L-1}}^{b^L-1}
    f(m)
    \log_b\!\left(\frac{m+1}{m}\right)
\]
admit exact product, Gamma-function, or asymptotic evaluations analogous
to those obtained here for \(s_b\).

Such extensions would connect the present Benford-weighted calculation
more directly with the classical theory of digital and \(q\)-additive
functions.  We do not pursue them here.


\subsection{Summary of the boundary correction}
\label{subsec:boundary-correction-summary}

The principal phenomenon established in this paper is a summable
boundary effect.

The significant positions closest to the leading edge retain a
systematic logarithmic bias. Although the bias at position \(j\)
decays exponentially with \(j\), the sum of these deviations converges
to the finite constant
\[
    c_b.
\]
Consequently, the mean sum of the first \(J\) significant digits is not
merely
\[
    J\frac{b-1}{2}+o(1),
\]
but rather
\[
    J\frac{b-1}{2}-c_b+o(1)
\]
whenever the underlying orbit is sampled with sufficient quantitative
equidistribution.

This provides a precise sense in which the logarithmic structure of
the leading digits survives beyond the first-order uniform-digit law.
The extent to which analogous second-order effects persist deeper into
the full digital expansion remains open.

% ============================================================
% Section 8
% ============================================================

\section{Conclusion}
\label{sec:conclusion}

This paper studied the cumulative effect of the nonuniform early
significant digits under the logarithmic, or Benford, mantissa measure.
Starting from the classical distribution of finite significant
prefixes, we obtained an exact logarithmically weighted finite
digit-sum identity and used it to evaluate the cumulative signed
significant-digit bias in closed form:
\[
\boxed{
    c_b
    =
    \frac b2
    -
    \frac{b-1}{2}\log_b(2\pi).
}
\]
The same finite identity also yields explicit control of the approach
to this limit through the Stirling expansion and its positive-real
remainder bound.

The second part of the paper transferred this model-side constant to
deterministic exponential sequences.  For multiplicatively independent
integers \(a,b\ge2\), quantitative equidistribution of
\(\{n\log_b a\}\) allows a logarithmically growing block of leading
significant digits to be averaged with absolute second-order accuracy.
For an admissible growing window \(J_N\), this gives
\[
\boxed{
    \mathcal S_{N,J_N}(a;b)
    =
    J_N\frac{b-1}{2}
    -
    c_b
    +
    o(1).
}
\]
Thus the bounded correction predicted by the logarithmic mantissa
model survives in an actual deterministic orbit.

The first-order ingredients behind this result belong to established
theory: finite significant-digit laws, higher-digit Gamma
representations, uniform distribution modulo one, and discrepancy
methods for exponential sequences all have substantial prior
literature.  The contribution developed here is instead the finite
Benford-weighted digit-sum evaluation, its closed cumulative constant,
and the retention of that same constant as a second-order term in the
controlled deterministic leading window.

The present method does not address the complete digit expansion of
\(a^n\), for which the number of digits grows linearly with \(n\),
nor does it establish variance laws, independence, mixing, or
probabilistic limit theorems for full digit sums.  Extending
second-order control beyond logarithmically growing leading windows,
and understanding whether related boundary corrections occur for more
general digital functions, remain natural problems for future work.

\section*{Code and reproducibility}

The computational verification materials described in
Section~\ref{sec:verification} are archived as version 1.0.0 on Zenodo
\cite{Dhoundiyal2026Reproducibility}. These computations are
supplementary to the analytic proofs and are not required for the
validity of the main results.

\section*{Acknowledgments}

The author used AI-assisted tools, including ChatGPT and Gemini Deep
Research, for literature-search assistance, adversarial review,
computational-verification support, and editorial refinement.
AI-generated suggestions were not treated as mathematical evidence,
and the author takes full responsibility for all mathematical claims,
proofs, citations, and conclusions in the manuscript.

\bibliographystyle{plain}
\bibliography{references4}

\end{document}
